\section{Related Work}
\label{sec:related}

\paragraph{Neural Cellular Automata.}
Neural Cellular Automata learn a local update rule that, applied
repeatedly to a grid of hidden states, gives rise to globally
coordinated behavior~\cite{mordvintsev2020growing}. Subsequent work has
demonstrated NCAs that grow morphologies from a seed, classify under
distribution shift, perform texture synthesis, and self-organize
attention~\cite{niklasson2021selforganising,palm2022variational,najarro2022hypernca}.
NCAs are typically trained with stochastic update schedules and a
``pool'' of intermediate states for stability. Recent work treats the
NCA's iteration count as a meaningful axis of computation and studies
when more iterations buy more reasoning. Our use of NCAs as
\emph{action-conditioned world models} extends this line: rather than
asking the NCA to converge to a static target, we ask it to predict
the engine's deterministic next state at every tick.

\paragraph{World models and self-prediction.}
World-model literature dates back at least to Schmidhuber's curiosity
work and was popularized for vision-based RL by
Ha and Schmidhuber~\cite{ha2018world}. Recent transformer-based
world models~\cite{micheli2022iris,robine2023transformer} compress
images to discrete tokens and roll out dynamics autoregressively over
those tokens; PWM-style methods integrate planning into a learned
model. Our problem domain---fully discrete grid observations,
authored rule sets, dozens of related games sharing a substrate---is
closer to procedural-content benchmarks like NetHack and CHAI than to
Atari. Compared to MuZero-style models that learn an abstract latent
dynamics, our model commits to predicting the engine's next state in
its native grid representation, which makes ground-truth supervision
exact and lets us compare to the engine cell-by-cell.

\paragraph{PuzzleScript and rule-based puzzle games.}
PuzzleScript~\cite{lavelle2013puzzlescript} is a domain-specific
language for grid-based puzzle games that supports a rich family of
authoring patterns: chained pushing (\texttt{[ Crate | Crate ] -> [
Crate | Crate ]}), \texttt{again}-driven simulation (gravity, ball
trajectories, light beams), multi-character control, and rule
priorities. Several hundred community games are publicly available and
share both a tokenization scheme (the engine source) and a common
state encoding, making the family a natural multi-game benchmark for
world models. Prior work has studied PuzzleScript primarily for
procedural content generation and human play; to our knowledge, ours
is the first systematic study of learned forward models on a broad
PuzzleScript gallery.

\paragraph{Adaptive computation.}
PonderNet~\cite{banino2021pondernet} learns a per-input halting
distribution over a recurrent module, with a KL prior over the number
of iterations. Adaptive Computation Time~\cite{graves2016act} and
Universal Transformers~\cite{dehghani2018universal} are predecessors.
Our adaptive-halt experiments adapt PonderNet to NCA world models
with a shared body, so that ``halt at step $k$'' coherently asks the
same rule set to converge in $k$ iterations rather than selecting
between $k$ different models.

\paragraph{Learned local-update grids more broadly.}
ConvLSTMs~\cite{shi2015convlstm}, graph neural networks, and
diffusion-style iterative refinement all share with NCAs the high-level
idea of repeated local computation. NCAs differ in committing to a
\emph{shared} (across iterations) update rule on a fixed grid, which
matches the PuzzleScript engine's outer-loop semantics. A detailed
comparison to procedural / engine-aligned NCA work is deferred to
the appendix.
